\documentclass[11pt,a4paper]{article}
\usepackage[utf8]{inputenc}
\usepackage[T1]{fontenc}
\usepackage{lmodern}
\usepackage{microtype}
\usepackage{booktabs}
\usepackage{graphicx}
\usepackage{amsmath}
\usepackage{amssymb}
\usepackage{siunitx}
\usepackage[hidelinks]{hyperref}
\usepackage[margin=1in]{geometry}
\usepackage{authblk}
\usepackage{xcolor}
\usepackage{natbib}
\bibliographystyle{plainnat}

\title{A SUMO E3 ligase promotes long non-coding RNA transcription to regulate small RNA-directed DNA elimination}
\author[1]{Anonymous Author}
\affil[1]{Institute of Molecular Biology}
\date{}

\begin{document}
\maketitle

\begin{abstract}
Small RNAs target their complementary chromatin regions for gene silencing through nascent long non-coding RNAs (lncRNAs). In the ciliated protozoan \textit{Tetrahymena thermophila}, Piwi-associated small RNAs (scnRNAs) interact with nascent lncRNA transcripts derived from the somatic macronucleus (MAC) genome to drive target-directed small RNA degradation (TDSD); the residual scnRNAs that escape TDSD subsequently direct programmed DNA elimination of germline-limited internal eliminated sequences (IESs). Here we show that the SUMO E3 ligase Ema2 is required for the accumulation of MAC-derived lncRNAs and, consequently, for TDSD and for completion of DNA elimination that is necessary to generate viable sexual progeny. Ema2 directly interacts with the SUMO E2 conjugating enzyme Ubc9 and enhances SUMOylation of the transcription elongation regulator Spt6. Loss of Ema2 reduces the high molecular weight SUMO-conjugate population (mainly $>$200~kDa) to $\sim$50\% of wild-type levels at conjugation stages, and it compromises the association of Spt6 and RNA polymerase II with chromatin. Together, these results indicate that Ema2-directed SUMOylation actively promotes lncRNA transcription, a prerequisite for communication between the somatic genome and small RNAs during programmed genome rearrangement.
\end{abstract}

\section{Introduction}

Small non-coding RNAs guide sequence-specific modifications of chromatin across eukaryotes, ranging from piRNA-directed transposon silencing in animal germlines to the RNA-induced transcriptional silencing pathway in fission yeast \citep{Grewal2010,Czech2018}. A common feature of these pathways is that small RNAs do not engage DNA directly but instead base-pair with nascent long non-coding RNA (lncRNA) transcripts produced by RNA polymerase II (Pol II) at the target loci \citep{Moazed2009,Holoch2015}. The resulting small-RNA/nascent-lncRNA interaction recruits effector complexes that deposit repressive chromatin marks or, in some contexts, trigger degradation of the guiding small RNAs themselves. The availability of chromatin-associated lncRNAs is therefore a rate-limiting step for small-RNA-directed genome regulation.

The ciliated protozoan \textit{Tetrahymena thermophila} provides a powerful system to dissect this logic because it separates germline and somatic functions into two nuclei within the same cell: a transcriptionally silent germline micronucleus (MIC) and a transcriptionally active somatic macronucleus (MAC) \citep{Karrer2012}. During the sexual cycle (conjugation), a new MAC is differentiated from a zygotic nucleus, and approximately 12{,}000 internal eliminated sequences (IESs) are excised from the developing MAC in a process called programmed DNA elimination \citep{Chalker2013}. This process is directed by $\sim$26--32~nt Piwi-associated small RNAs called scan RNAs (scnRNAs), which are produced from the MIC, imported into the parental MAC, and subjected to target-directed small RNA degradation (TDSD) upon base-pairing with nascent lncRNAs transcribed from the parental MAC genome \citep{Mochizuki2002,Aronica2008}. ScnRNAs not consumed by TDSD are retained and subsequently guide elimination of the complementary germline-limited sequences in the new MAC.

This model makes a strong prediction: factors that promote the production of parental-MAC lncRNAs should be essential for TDSD and, downstream of TDSD, for IES elimination and sexual progeny viability. Several chromatin and transcription elongation factors have been implicated in this step, including the histone chaperone and Pol II-associated elongation factor Spt6 \citep{Hartzog2013}, but the upstream regulators that license lncRNA transcription from the parental MAC remain poorly understood. Post-translational modification by the small ubiquitin-like modifier (SUMO) is an attractive candidate, because SUMOylation is known to control transcription elongation factor function and chromatin association in other systems \citep{Hendriks2016,Geiss2007}, and because SUMO E3 ligases provide target specificity to an otherwise promiscuous SUMO E2 conjugating enzyme, Ubc9 \citep{Gareau2010}.

In this study, we identify Ema2 as a \textit{Tetrahymena} SUMO E3 ligase that is specifically required for scnRNA-directed DNA elimination. Using genetic deletion of \textit{EMA2}, in vitro biochemistry, SUMO-conjugate profiling, cell fractionation, and targeted analysis of parental-MAC lncRNAs, we show that Ema2 binds the SUMO E2 Ubc9, enhances SUMOylation of Spt6, promotes Spt6 and Pol II association with chromatin, sustains the accumulation of lncRNAs from IES-containing MAC loci, and is required for TDSD, IES elimination, and progeny viability. These findings place a SUMO E3 ligase directly upstream of the lncRNA step that couples the somatic genome to the scnRNA surveillance machinery.

\section{Methods}

\subsection{Strains, growth, and conjugation}
\textit{Tetrahymena thermophila} strains, including wild-type (CU428) and \textit{EMA2} knockout (\textit{EMA2} KO) strains, were grown in SPP medium containing 2\% proteose peptone at 30~$^\circ$C \citep{Gorovsky1975}. For conjugation, growing cells at $\sim$5--7~$\times$10$^{5}$ cells/mL of two different mating types were washed, pre-starved for $\sim$12--24~h, and mixed in 10~mM Tris (pH~7.5) at 30~$^\circ$C. The time of mixing was defined as 0~h post-mixing (hpm). Cells were incubated in 10~mM Tris pH~7.5 at 30~$^\circ$C until 24~hpm and then fed by adding an equal volume of 2$\times$ SPP.

\subsection{Viability test for sexual progeny}
To test progeny viability, mating pairs were isolated into drops of 1$\times$ SPP at 6~hpm and incubated for $\sim$3 days. Completion of conjugation in wild-type and \textit{EMA2} KO pairs was further assessed by 6-methylpurine (6-mp) resistance and paromomycin sensitivity as previously described \citep{Mochizuki2002}.

\subsection{Small RNA analyses}
For northern hybridization of scnRNAs, 10~$\mu$g of total RNA isolated from conjugating wild-type and \textit{EMA2} KO cells using TRIzol reagent was separated on a 15\% acrylamide-urea gel and analyzed by northern blot using a radio-labeled Mi-9 probe as previously described \citep{Aronica2008}. For small RNA sequencing, normalized numbers (RPKM [reads per kilobase of unique sequences per million]) of 26- to 32-nt small RNAs that uniquely matched MDS and IES tiles were counted. Because DNA elimination is completely blocked in \textit{TWI1} KO cells, a correction factor of 0.35 was applied to set the average IES retention index of \textit{TWI1} KO cells to 1.

\subsection{Recombinant protein production and GST pull-down}
\textit{E.~coli} codon-optimized synthetic \textit{EMA2} and \textit{UBC9} genes (\textit{EMA2-Ec} and \textit{UBC9-Ec}) were cloned into the BamHI/XhoI sites of pGEX-4T-1 and the NdeI/BamHI sites of pET28a(+), respectively, to obtain pGEX-4T-1-EMA2-Ec and pET28-UBC9-Ec. pET28-UBC9-Ec was transformed into \textit{E.~coli} BL21. His-Ubc9 expression was induced for $\sim$16~h at 25~$^\circ$C with 0.05~mM IPTG. Cells were lysed in HEPES buffer (50~mM HEPES pH~7.5, 800~mM KCl, 10\% glycerol, 0.2~mM PMSF, 1$\times$ cOmplete EDTA-free [Roche]) containing 10~mM imidazole, incubated with Ni-NTA beads for 2~h at 4~$^\circ$C, washed four times with HEPES buffer containing 20~mM imidazole, eluted with HEPES buffer containing 250~mM imidazole, dialyzed into 50~mM Tris pH~7.5, 150~mM NaCl, 2~mM MgCl$_2$, 1~mM DTT, 50\% glycerol, and passed through a Microcon 10~kDa filter (Merck).

For GST pull-down, BL21 \textit{E.~coli} cells were transformed with pGEX-4T-1 or pGEX-4T-1-EMA2-Ec, and expression of GST-alone or GST-Ema2 was induced for $\sim$16~h at 25~$^\circ$C. Cells were lysed in 50~mM Tris-HCl pH~7.5, 100~mM NaCl, 1~mM DTT, 0.2~mM PMSF, 1$\times$ cOmplete EDTA-free and incubated with Glutathione Sepharose 4B for 90~min at 4~$^\circ$C. Beads were then incubated with 10~$\mu$g of His-Ubc9 in GST pull-down buffer (20~mM Tris-HCl pH~7.5, 100~mM NaCl, 1\% NP-40, 5~mM MgCl$_2$, 1~mM DTT, 0.2~mM PMSF, 1$\times$ cOmplete EDTA-free) for 90~min at 4~$^\circ$C, washed four times with GST pull-down buffer, and eluted with 2$\times$ SDS sample buffer.

\subsection{HA-Smt3 and His-Smt3 expressing strains and SUMO-conjugate purification}
\textit{SMT3} cDNA was amplified by RT-PCR from total RNA of vegetatively growing CU428 using SMT3\_HA\_FW and SMT3\_HisExt\_RV and cloned into the BamHI/SpeI sites of pBNMB1-HA using NEBuilder HiFi DNA assembly to obtain pBNMB1-HA-SMT3. Transformed cells were selected in 300~$\mu$g/mL puromycin and 1~$\mu$g/mL CdCl$_2$ and phenotypically assorted as previously described \citep{Hamilton2000} until cells grew in 1.2~mg/mL puromycin.

After 3~h at 30~$^\circ$C, neo (and His-Smt3 from pBCMB1-His-SMT3) expression was induced by adding 1~$\mu$g/mL CdCl$_2$ and cells were incubated for an additional 1~h. For purification of His-Smt3 conjugates, cells were lysed in Guanidine Lysis Buffer (6~M guanidine-HCl, 93.2~mM Na$_2$HPO$_4$, 6.8~mM NaH$_2$PO$_4$, 10~mM Tris-HCl pH~8.0) by sonication (2~$\times$~25 pulses, 20\% power, 60\% duty cycle) using an Omni-Ruptor 250 Ultrasonic Homogenizer. After centrifugation, the cleared lysate was supplemented with 50~mM imidazole and 5~mM $\beta$-mercaptoethanol and incubated with Ni-NTA agarose beads. Beads were washed and eluted in Elution buffer (6~M urea, 58~mM Na$_2$HPO$_4$, 42~mM NaH$_2$PO$_4$, 10~mM Tris-HCl pH~8.0, 500~mM imidazole). Eluted proteins were precipitated in 10\% TCA, resuspended in 1$\times$ SDS buffer, and analyzed by western blot and mass spectrometry.

For HA-Smt3 conjugate purification, total proteins were precipitated by incubating cells in 6\% trichloroacetic acid (TCA) on ice for 10~min, spun down at 13{,}000~rpm for 5~min at 4~$^\circ$C, dissolved in 1$\times$ SDS sample buffer, neutralized with 2~M Tris, and incubated for 5~min at 95~$^\circ$C. The dissolved proteins (600~$\mu$L) were suspended in 7.4~mL of 20~mM Tris pH~7.5, 100~mM NaCl, 2~mM MgCl$_2$, 2~mM CaCl$_2$ and incubated with EZview Red anti-HA agarose beads (Sigma) overnight at 4~$^\circ$C. Beads were washed 5 times with 20~mM Tris pH~7.5, 100~mM NaCl, 2~mM MgCl$_2$, 2~mM CaCl$_2$, 0.1\% Tween~20, and bound proteins were eluted with 1$\times$ SDS buffer at 95~$^\circ$C for 5~min.

\subsection{RT-PCR analyses of lncRNA transcripts}
Total RNA from mating wild-type and \textit{EMA2} KO cells was isolated at 6~hpm using TRIzol (Invitrogen). lncRNA transcripts from the MAC loci containing the M-, L8-, and R2 IES boundaries were amplified by nested PCR (95~$^\circ$C, 20~s / 50~$^\circ$C, 30~s / 68~$^\circ$C, 1~min; 30 cycles each) with primers M5$^\prime$-3/M3$^\prime$-3 (M-1st), M5$^\prime$-4/M5$^\prime$-4 (M-2nd), L8\_5$^\prime$-1/L8\_3$^\prime$-1 (L8-1st), L8\_5$^\prime$-2/L8\_3$^\prime$-2 (L8-2nd), R2\_5$^\prime$-1/R2\_3$^\prime$-1 (R2-1st), and R2\_5$^\prime$-2/R2\_3$^\prime$-2 (R2-2nd), and products were analyzed on 1\% agarose gels. The constitutively expressed \textit{RPL21} mRNA detected with RPL21-FW/RPL21-RV served as a positive control.

\subsection{Cell fractionation}
Cell fractionation was performed as described \citep{Ali2018}. $1.4 \times 10^{6}$ cells were incubated with 1~mL of Lysis buffer (50~mM Tris pH~8.0, 100~mM NaCl, 5~mM MgCl$_2$, 1~mM EDTA, 0.5\% Triton X-100, 1$\times$ Complete Protease Inhibitor) on ice for 30~min and spun at 15{,}000$\,g$ for 10~min at 4~$^\circ$C. Proteins in 500~$\mu$L of the supernatant (soluble fraction) were precipitated with 10\% (final) TCA, washed three times with ice-cold acetone, and dissolved in 70~$\mu$L of 1$\times$ SDS sample buffer. The insoluble pellet was resuspended in 2~mL fresh Lysis buffer, sonicated (5 cycles, 20\% power, 50\% duty cycle, 5 pulses, with $>$20~s intervals) using an Omni-Ruptor 250, and centrifuged at 20{,}000$\,g$ for 10~min at 4~$^\circ$C; proteins in 1~mL of the resulting supernatant (solubilized chromatin) were TCA-precipitated and dissolved in 70~$\mu$L 1$\times$ SDS buffer. For RNase treatment, lysis was performed as above but with 20~$\mu$g/mL RNase~A, and the final insoluble pellet was directly dissolved in 140~$\mu$L of 1$\times$ SDS sample buffer.

\subsection{\textit{SPT6} germline knockout}
To create a knockout construct for \textit{SPT6}, the 5$^\prime$ and 3$^\prime$ flanking regions of the \textit{SPT6} locus were amplified by PCR with SPT6-KO-5F-FW/SPT6-KO-5F-RV and SPT6-KO-3F-FW/SPT6-KO-3F-RV, respectively, and concatenated by overlap PCR using SPT6-KO-5F-FW and SPT6-KO-3F-RV.

\section{Results}

\subsection{Ema2 is a SUMO E3 ligase that binds the E2 conjugating enzyme Ubc9}
A functional SUMO E3 ligase must bind Ubc9 to orient SUMO transfer onto substrates \citep{Gareau2010}. To test whether Ema2 shares this defining biochemical feature, we purified recombinant GST-Ema2 from \textit{E.~coli} and incubated glutathione-Sepharose-bound GST or GST-Ema2 with purified His-Ubc9. After four washes, bound His-Ubc9 was eluted in 2$\times$ SDS sample buffer and detected by western blot. His-Ubc9 was specifically retained on GST-Ema2 beads but not on GST-alone beads, identifying Ema2 as a direct Ubc9 interaction partner and supporting its classification as a SUMO E3 ligase.

\subsection{Ema2 is required for the accumulation of high-molecular-weight SUMO conjugates during conjugation}
To ask whether Ema2 affects the global SUMO conjugation profile during the sexual cycle, we compared His-Smt3 conjugates purified under denaturing conditions from wild-type and \textit{EMA2} KO conjugating cells at two conjugation time points. Both time points displayed a characteristic ladder of high molecular weight SUMO conjugates, predominantly $>$200~kDa, in the wild-type cross. These high-MW conjugates were strongly reduced in the \textit{EMA2} KO cross, dropping to approximately 50\% of wild-type levels at both time points (Table~\ref{tab:sumo}). Thus, Ema2 is required for the proper build-up of a conjugation-specific SUMO sub-proteome in \textit{Tetrahymena}.

\begin{table}[h]
\centering
\caption{SUMO-conjugate profile in wild-type versus \textit{EMA2} KO conjugating cells. Values reflect semi-quantitative analysis of His-Smt3 conjugates above 200~kDa (see Methods).}
\label{tab:sumo}
\begin{tabular}{lcc}
\toprule
Genotype & High-MW SUMO conjugate level & Dominant MW range \\
\midrule
Wild type (WT)   & $\sim$100\% (reference)        & mainly $>$200~kDa \\
\textit{EMA2} KO & $\sim$50\% of WT               & mainly $>$200~kDa \\
\bottomrule
\end{tabular}
\end{table}

\subsection{Ema2 promotes Spt6 SUMOylation and chromatin association of Spt6 and Pol II}
Mass spectrometry of HA-Smt3 pulldowns followed by targeted western blot identified the transcription elongation factor Spt6 as a SUMO substrate whose modification was diminished in \textit{EMA2} KO cells, consistent with Ema2 acting as the cognate E3 ligase for Spt6 SUMOylation during conjugation. Because SUMOylation of elongation factors often controls their recruitment to chromatin in other systems \citep{Hendriks2016,Geiss2007}, we fractionated wild-type and \textit{EMA2} KO cells into soluble and solubilized-chromatin pools using the conditions summarized in Table~\ref{tab:frac}. Both Spt6 and RNA polymerase II were enriched in the chromatin fraction of wild-type cells but were depleted from the chromatin fraction of \textit{EMA2} KO cells, demonstrating that Ema2 is required for the stable association of the elongation machinery with chromatin during conjugation.

\begin{table}[h]
\centering
\caption{Key parameters of the cell fractionation assay used to monitor chromatin association of Spt6 and Pol II.}
\label{tab:frac}
\begin{tabular}{lcc}
\toprule
Parameter & Value & Notes \\
\midrule
Input cells                       & $1.4\times10^{6}$ & per fractionation \\
Lysis buffer Triton X-100         & 0.5\%             & on ice, 30 min \\
Soluble-fraction spin             & $15{,}000\,g$, 10 min, 4~$^\circ$C & \\
Chromatin-fraction spin           & $20{,}000\,g$, 10 min, 4~$^\circ$C & after sonication \\
Sonication (chromatin release)    & 5 cycles          & 20\% power, 50\% duty \\
RNase A (where indicated)         & 20~$\mu$g/mL      & included in lysis \\
\bottomrule
\end{tabular}
\end{table}

\subsection{Ema2 is required for accumulation of parental-MAC lncRNAs at IES-boundary loci}
To directly test whether the compromised Pol II--chromatin association in \textit{EMA2} KO cells translates into a loss of non-coding transcription, we isolated total RNA from mating wild-type and \textit{EMA2} KO cells at 6~hpm and amplified lncRNA transcripts spanning the M, L8, and R2 IES-boundary regions by nested PCR (Table~\ref{tab:lncrna}). lncRNA products from all three IES-boundary loci were readily detected in wild-type mating cells but were markedly reduced or absent in \textit{EMA2} KO cells, while the constitutively expressed \textit{RPL21} mRNA served as a loading and RT efficiency control and was detected at comparable levels in both genotypes. Thus, Ema2 supports the accumulation of parental-MAC lncRNAs that must base-pair with scnRNAs for TDSD.

\begin{table}[h]
\centering
\caption{Targeted loci and primer sets used for nested RT-PCR of parental-MAC lncRNAs at IES boundaries.}
\label{tab:lncrna}
\begin{tabular}{llll}
\toprule
Locus & First-round primers & Second-round primers & Control \\
\midrule
M-boundary   & M5$^\prime$-3 + M3$^\prime$-3     & M5$^\prime$-4 + M5$^\prime$-4     & \textit{RPL21} mRNA \\
L8-boundary  & L8\_5$^\prime$-1 + L8\_3$^\prime$-1 & L8\_5$^\prime$-2 + L8\_3$^\prime$-2 & \textit{RPL21} mRNA \\
R2-boundary  & R2\_5$^\prime$-1 + R2\_3$^\prime$-1 & R2\_5$^\prime$-2 + R2\_3$^\prime$-2 & \textit{RPL21} mRNA \\
\bottomrule
\end{tabular}
\end{table}

All PCR reactions were run for 30 cycles (95~$^\circ$C, 20~s / 50~$^\circ$C, 30~s / 68~$^\circ$C, 1~min) and resolved on 1\% agarose gels.

\subsection{Ema2 is required for TDSD and for viable sexual progeny}
Loss of parental-MAC lncRNAs predicts failure of TDSD and, downstream, of IES elimination. We monitored the scnRNA population in conjugating wild-type and \textit{EMA2} KO cells by northern blot using the Mi-9 probe on 10~$\mu$g of total RNA separated on a 15\% acrylamide-urea gel. In wild-type cells, scnRNAs progressively declined over the course of conjugation, consistent with active TDSD; in \textit{EMA2} KO cells, scnRNAs persisted, indicating defective TDSD. Consistent with a failure of scnRNA-directed DNA elimination, small-RNA sequencing of 26--32~nt reads mapped to MDS and IES tiles (expressed as RPKM) revealed a shift in the MDS/IES distribution, and a correction factor of 0.35---empirically defined in \textit{TWI1} KO cells to set their IES retention index to 1---was used to normalize IES retention indices across genotypes. Finally, we isolated mating pairs of wild-type and \textit{EMA2} KO cells into individual drops of 1$\times$ SPP at 6~hpm and scored sexual progeny viability $\sim$3 days later. Completion of conjugation was assessed using 6-methylpurine (6-mp) resistance and paromomycin sensitivity as orthogonal markers, as described \citep{Mochizuki2002}. \textit{EMA2} KO cells produced substantially fewer viable progeny than wild-type, establishing that Ema2 is required for productive sexual reproduction (Table~\ref{tab:phen}).

\begin{table}[h]
\centering
\caption{Phenotypic summary of \textit{EMA2} KO cells in the scnRNA/DNA elimination pathway.}
\label{tab:phen}
\begin{tabular}{lcc}
\toprule
Readout & Wild type & \textit{EMA2} KO \\
\midrule
Progeny viability (6-mp/paromomycin assay)         & normal            & reduced \\
High-MW SUMO conjugates ($>$200~kDa)               & reference         & $\sim$50\% of WT \\
Spt6 SUMOylation                                   & present           & reduced \\
Spt6 / Pol II chromatin fraction                   & enriched          & depleted \\
Parental-MAC lncRNA (M, L8, R2 boundaries)         & detected          & reduced/absent \\
scnRNA TDSD                                        & normal decline    & persistent scnRNAs \\
\bottomrule
\end{tabular}
\end{table}

\section{Discussion}

The central finding of this study is that a SUMO E3 ligase, Ema2, sits upstream of a non-coding transcription step that is essential for small-RNA-directed DNA elimination in \textit{Tetrahymena}. By biochemistry, Ema2 binds the SUMO E2 conjugating enzyme Ubc9 and enhances SUMOylation of the elongation factor Spt6. By cell biology, Ema2 is required for the chromatin-association of Spt6 and RNA polymerase II during conjugation and for the accumulation of high-molecular-weight SUMO conjugates. By genetics, loss of \textit{EMA2} collapses parental-MAC lncRNAs at IES-boundary loci, impairs TDSD of scnRNAs, compromises programmed DNA elimination, and reduces sexual progeny viability. These observations link a specific SUMOylation event to the lncRNA output that small RNAs need in order to distinguish somatic from germline sequences, providing a molecular explanation for why SUMOylation is required in this pathway.

Our results fit naturally with a broader view in which SUMOylation acts as a ``solubility and recruitment'' switch for transcription elongation factors \citep{Hendriks2016,Geiss2007}. In other systems, SUMOylation has been proposed to stabilize multi-protein elongation complexes on chromatin and to fine-tune the transition between initiation and productive elongation. Our data extend this logic to a developmentally induced non-coding transcription program: during conjugation, when the parental MAC must transiently produce lncRNAs throughout its genome to guide scnRNA scanning, Ema2-directed SUMOylation of Spt6 appears to be the permissive signal that licenses Pol II to traverse chromatin. The observation that both Spt6 and Pol II are simultaneously depleted from the chromatin fraction in \textit{EMA2} KO cells is consistent with this model and with the known role of Spt6 in Pol II progression through chromatinized templates \citep{Hartzog2013}.

The reduction of high-MW SUMO conjugates to $\sim$50\% in \textit{EMA2} KO conjugating cells is notable because it argues that Ema2 contributes a quantitatively large, but not exclusive, share of the conjugation-stage SUMO signal. Residual SUMOylation in the \textit{EMA2} KO cross is consistent with the existence of additional SUMO E3 ligases or with E2-only conjugation onto certain substrates, and it implies that the Ema2-dependent sub-proteome is enriched for transcription-linked targets. Identifying those targets systematically---for example, by combining His-Smt3 pulldown with quantitative mass spectrometry across a time course of conjugation---will be a natural next step and should reveal whether Ema2 SUMOylates additional elongation or chromatin factors alongside Spt6.

Our findings also have implications beyond \textit{Tetrahymena}. The requirement that small-RNA-guided genome modification passes through a nascent-lncRNA intermediate is conserved from fission yeast heterochromatin to mammalian piRNA pathways \citep{Moazed2009,Holoch2015,Czech2018}. In each of these systems, any factor that licenses chromatin-associated lncRNA output would be predicted to have a gate-keeping role analogous to that of Ema2. The identification of a SUMO E3 ligase as such a factor in \textit{Tetrahymena} therefore raises the possibility that SUMO-dependent elongation control is a conserved upstream node in small-RNA-directed genome regulation.

Several limitations should be noted. First, the reduction of high-MW SUMO conjugates and of chromatin-bound Spt6/Pol II in \textit{EMA2} KO cells is measured at the level of a conjugation-stage population and does not resolve the precise kinetics of the defect. Second, although our data establish that Ema2 binds Ubc9 and that Spt6 SUMOylation is Ema2-dependent, we have not yet reconstituted Ema2-stimulated Spt6 SUMOylation in a fully defined system; such an assay, ideally in combination with SIM-motif mutations in Spt6, would formally establish direct E3 activity on Spt6. Third, the lncRNA RT-PCR assay is targeted to three IES-boundary loci and is not genome-wide; a strand-specific nascent-RNA sequencing approach would be required to quantify the full extent of the lncRNA defect and to ask whether it affects all IES-adjacent loci uniformly. Finally, the \textit{SPT6} germline knockout reported here provides the genetic tool to directly test the prediction that Spt6 itself is required for lncRNA production and TDSD, closing the loop between SUMO substrate and pathway phenotype.

Together, these findings define Ema2 as a SUMO E3 ligase whose activity licenses the lncRNA transcription that small RNAs need in order to target the right chromatin regions, and they establish SUMO-dependent elongation control as a core regulatory layer of programmed DNA elimination in \textit{Tetrahymena}.

\bibliography{references}

\end{document}
